We consider a reasoning policy interacting with a prefix \(p_t\) and a remaining token budget \(b_t\).
At each decision point the controller chooses one of three actions:
\[
\mathcal{A} = \{\continue, \revise, \abstain\}.
\]
Here \continue means extending the current trajectory, \revise means taking one bounded corrective intervention, and \abstain means stopping and declining further computation.

\paragraph{Oracle utility.}
For a prefix \(p_t\) in an exact-checker domain, we estimate an action utility
\[
\mathcal{U}(a \mid p_t, b_t)
\]
by rollout under a fixed base reasoner and exact task checker.
The oracle action is the action with maximal expected utility.
Because the environment is verifiable, the oracle can be computed at the prefix level rather than inferred only from final-task success.
In the experiments reported here, \abstain has fixed utility \(0\), while \continue and \revise are estimated from \(4\)--\(8\) Monte Carlo rollouts of a fixed base reasoner \(\pi_0\).
Each rollout is scored by
\[
\mathbb{1}\{\text{success}\} - \lambda_{\text{tok}} \cdot \text{tokens} - \gamma_{\text{wrong}} \cdot \mathbb{1}\{\text{unsafe wrong}\},
\]
so the benchmark trades off correctness, token cost, and harmful incorrect finalization.
A prefix is marked \emph{ambiguous} when the top-two estimated actions are not cleanly separated: either the top-two mean-utility gap is below \(0.05\), or their \(95\%\) confidence intervals overlap.
Main controller comparisons exclude ambiguous prefixes by default.

\paragraph{Ordered scalar controllers.}
The ordered-scalar controller class reduces the decision to a scalar score \(s(p_t)\) together with budget-dependent thresholds.
Abstractly, it has the form
\[
\pi_{\text{scalar}}(p_t, b_t) = g(s(p_t), b_t),
\]
where \(g\) is an order-preserving threshold policy over \(s\).
This covers natural scalar surrogates together with learned scalar summaries paired with order-preserving decision rules.
At the same time, the paper does \emph{not} claim to rule out every one-dimensional encoding.
The learned one-dimensional controller in the main text should be read as a strong compression baseline, not as the formal impossibility target.
Our benchmark is designed to test whether such scalar-ordering views are expressive enough for low-regret prefix control.

\paragraph{Structured controllers.}
We also study a factorized controller family that predicts a structured \processstate and then maps that state to action values.
The exact semantic content of the state is intentionally narrow.
One component captures local invalidity risk \(q_t\).
A second component captures \emph{continuation outcome statistics}: summary information about what happens if the current prefix is allowed to continue.
We avoid calling these richer features the full \(S_t\) from the original proposal, because in the current exact-checker domains the original Beta-style success uncertainty often degenerates.

\paragraph{Controller families.}
The main-text comparison includes five controller types:
\begin{enumerate}[leftmargin=1.5em]
  \item \orderedscalar: an ordered threshold controller over \(\mu_{\continue}\),
  \item \learnedoned: a learned one-dimensional ordered controller,
  \item \directpolicy: an unstructured direct policy,
  \item \triver exact-state: a structured controller with oracle state access,
  \item \triver predicted-state: one selected deployable predicted-state controller per domain.
\end{enumerate}
The exact-state controller is an upper bound on what the factorization can achieve when state is not the bottleneck.
The predicted-state controller measures deployment quality.
The predicted-state row is intentionally frozen for the main text.
Later appendix rescue probes are used to diagnose the deployment gap, not to reopen the core controller comparison.

\paragraph{Domains and metrics.}
We instantiate the benchmark on arithmetic and linear equations.
Our core metrics are oracle determinacy rate, crossing mass, action regret, oracle action accuracy, revision harm, and compute-value calibration.
For budgeted evaluation we also report action regret as a function of budget and equal-token frontiers.
